\subsection{\texorpdfstring{同构 \(\operatorname{Hom}_{\mathbf{B}}(\mathbf{E} \otimes_{\mathbf{A}} \mathbf{F}, \mathbf{G}) \to \operatorname{Hom}_{\mathbf{A}}(\mathbf{F}, \operatorname{Hom}_{\mathbf{B}}(\mathbf{E}, \mathbf{G}))\) 与 \(\operatorname{Hom}_{\mathbf{C}}(\mathbf{E} \otimes_{\mathbf{A}} \mathbf{F}, \mathbf{G}) \to \operatorname{Hom}_{\mathbf{A}}(\mathbf{E}, \operatorname{Hom}_{\mathbf{C}}(\mathbf{F}, \mathbf{G}))\)}{同构 \textbackslash operatorname\{Hom\}\_\{\textbackslash mathbf\{B\}\}(\textbackslash mathbf\{E\} \textbackslash otimes\_\{\textbackslash mathbf\{A\}\} \textbackslash mathbf\{F\}, \textbackslash mathbf\{G\}) \textbackslash 到 \textbackslash operatorname\{Hom\}\_\{\textbackslash mathbf\{A\}\}(\textbackslash mathbf\{F\}, \textbackslash operatorname\{Hom\}\_\{\textbackslash mathbf\{B\}\}(\textbackslash mathbf\{E\}, \textbackslash mathbf\{G\})) 且 \textbackslash operatorname\{Hom\}\_\{\textbackslash mathbf\{C\}\}(\textbackslash mathbf\{E\} \textbackslash otimes\_\{\textbackslash mathbf\{A\}\} \textbackslash mathbf\{F\}, \textbackslash mathbf\{G\}) \textbackslash 到 \textbackslash operatorname\{Hom\}\_\{\textbackslash mathbf\{A\}\}(\textbackslash mathbf\{E\}, \textbackslash operatorname\{Hom\}\_\{\textbackslash mathbf\{C\}\}(\textbackslash mathbf\{F\}, \textbackslash mathbf\{G\}))}}

设 E 为右 A-模，F 为左 A-模，G 为 Z-模，H 为由映射 \(f: E \times F \rightarrow G\) 组成的 Z-模，这些映射是 Z-双线性的且满足

\[
f (x \lambda , y) = f (x, \lambda y) \quad \text { for } x \in \mathrm{E}, y \in \mathrm{F}, \lambda \in \mathrm{A}.\tag{1}
\]

已经看到（§ 3，第 1 号，命题 1）存在一个典范的 Z-模同态

\[
\mathrm{H} \rightarrow \operatorname{Hom} _ {\mathbf {Z}} (\mathrm{E} \otimes_ {\mathrm{A}} \mathrm{F}, \mathrm{G}).\tag{2}
\]

另一方面，在 \(\mathrm{Hom}_{\mathbf{Z}}(\mathrm{E},\mathrm{G})\) 上已定义了左 A-模结构，在 \(\mathrm{Hom}_{\mathbf{Z}}(\mathrm{F},\mathrm{G})\) 上定义了右 A-模结构（§ 3，第 3 号）；因此我们可以考虑 Z-模 \(\mathrm{Hom}_{\mathbf{A}}(\mathrm{E},\mathrm{Hom}_{\mathbf{Z}}(\mathrm{F},\mathrm{G}))\) 和 \(\mathrm{Hom}_{\mathbf{A}}(\mathrm{F},\mathrm{Hom}_{\mathbf{Z}}(\mathrm{E},\mathrm{G}))\) 。一个映射 \(f\) （从 \(\mathbf{E} \times \mathbf{F}\) 到 \(\mathbf{G}\) 的）被典范地等同于一个从 \(\mathbf{E}\) 到集合 \(\mathbf{G}^{\mathbf{F}}\) （从 \(\mathbf{F}\) 到 \(\mathbf{G}\) 的映射所成的集合）中的映射（集合论，II，§ 5，第 2 号）；通过表达后一映射属于 \(\mathrm{Hom}_{\mathbf{A}}(\mathrm{E},\mathrm{Hom}_{\mathbf{Z}}(\mathrm{F},\mathrm{G}))\) 这一事实，我们恰好得到 \(f\) 是双加性的且满足条件 (1)；由此存在一个典范同构

\[
\mathrm{H} \rightarrow \operatorname{Hom} _ {\mathbf {A}} (\mathrm{E}, \operatorname{Hom} _ {\mathbf {Z}} (\mathrm{F}, \mathrm{G}))\tag{3}
\]

类似地，可定义一个典范同构

\[
\mathrm{H} \rightarrow \operatorname{Hom} _ {\mathbf {A}} (\mathrm{F}, \operatorname{Hom} _ {\mathbf {Z}} (\mathrm{E}, \mathrm{G})).\tag{4}
\]

现在假设 E 和 G 也具有左（相应地，右）B-模结构，并且 E 上的 A-模与 B-模结构是相容的。则 \(E \otimes_{A} F\) 典范地具有一个左（相应地，右）B-模结构（ \(\S 3\) ，第 4 号）；另一方面， \(\operatorname{Hom}_{\mathbf{B}}(\mathbf{E}, \mathbf{G})\) 具有典范的左 A-模结构（ \(\S 1\) ，第 14 号）。因此我们可以考虑Z-模 \(\operatorname{Hom}_{\mathbf{B}}(\mathbf{E} \otimes_{\mathbf{A}} \mathbf{F}, \mathbf{G})\) 和 \(\operatorname{Hom}_{\mathbf{A}}(\mathbf{F}, \operatorname{Hom}_{\mathbf{B}}(\mathbf{E}, \mathbf{G}))\) ，它们是 \(\operatorname{Hom}_{\mathbf{Z}}(\mathbf{E} \otimes_{\mathbf{A}} \mathbf{F}, \mathbf{G})\) 和 \(\operatorname{Hom}_{\mathbf{A}}(\mathbf{F}, \operatorname{Hom}_{\mathbf{Z}}(\mathbf{E}, \mathbf{G}))\) 的子模（ \(\S 2\) ，第1号，定理2）。我们考察在什么条件下，一个映射 \(f \in H\) 在同构（2）和（4）下的像分别是 \(\operatorname{Hom}_{\mathbf{B}}(\mathbf{E} \otimes_{\mathbf{A}} \mathbf{F}, \mathbf{G})\) 的一个元素和 \(\operatorname{Hom}_{\mathbf{A}}(\mathbf{E}, \operatorname{Hom}_{\mathbf{B}}(\mathbf{F}, \mathbf{G}))\) 的一个元素；在这两种情形中，我们都得到相同的条件

\[
\begin{array}{r l} f (\beta x, y) & = \beta f (x, y) \\ (\text { resp.   } f (x \beta , y) & = f (x, y) \beta) \end{array}
\]

对于 \(x \in \mathbf{E}, y \in \mathbf{F}, \beta \in \mathbf{B}\) 。

类似地，设F和G是左（相应地，右）C-模，并且F上的A-模与C-模结构是相容的。那么，一个映射 \(f \in H\) 在（2）或（3）下的像分别是 \(\operatorname{Hom}_{\mathbf{C}}(\mathbf{E} \otimes_{\mathbf{A}} \mathbf{F}, \mathbf{G})\) 或 \(\operatorname{Hom}_{\mathbf{A}}(\mathbf{E}, \operatorname{Hom}_{\mathbf{C}}(\mathbf{F}, \mathbf{G}))\) 的一个元素，当且仅当它满足相同的条件

\[
\begin{array}{c} f (x, \gamma y) = \gamma f (x, y) \\ (\text { resp. } f (x, y \gamma) = f (x, y) \gamma) \end{array}
\]

对于 \(x \in \mathbf{E}, y \in \mathbf{F}, \gamma \in \mathbf{C}\) 。

因此，我们已建立了如下结果（采用上文引入的记号）：

命题1. (a) 设E是一个(B, A)-双模，F是一个左A-模，G是一个左B-模。对于每一个映射 \(g \in \operatorname{Hom}_{\mathbf{B}}(\mathbf{E} \otimes_{\mathbf{A}} \mathbf{F}, \mathbf{G})\) ，设 \(g'\) 为从 F 到 \(\operatorname{Hom}_{\mathbf{B}}(\mathbf{E}, \mathbf{G})\) 中的映射，定义为 \((g'(y))(x) = g(x \otimes y)\) （对于 \(x \in \mathbf{E}, y \in \mathbf{F}\) ）。映射 \(g \mapsto g'\) 是一个同构

\[
\beta : \operatorname{Hom} _ {\mathbf {B}} (\mathbf {E} \otimes_ {\mathbf {A}} \mathbf {F}, \mathbf {G}) \rightarrow \operatorname{Hom} _ {\mathbf {A}} (\mathbf {F}, \operatorname{Hom} _ {\mathbf {B}} (\mathbf {E}, \mathbf {G})).\tag{5}
\]

\begin{enumerate}
\def\labelenumi{(\alph{enumi})}
\setcounter{enumi}{1}
\tightlist
\item
  设 \(\mathbf{E}\) 为一个右 \(\mathbf{A}\) -模， \(\mathbf{F}\) 一个 \((\mathbf{A}, \mathbf{C})\) -双模且 \(\mathbf{G}\) 一个右 \(\mathbf{C}\) -模。对于每个映射 \(h \in \operatorname{Hom}_{\mathbf{C}}(\mathbf{E} \otimes_{\mathbf{A}} \mathbf{F}, \mathbf{G})\) ，设 \(h'\) 是从 \(\mathbf{E}\) 到 \(\operatorname{Hom}_{\mathbf{C}}(\mathbf{F}, \mathbf{G})\) 中的映射，定义为 \((h'(x))(y) = h(x \otimes y)\) （对于 \(x \in \mathbf{E}\) , \(y \in \mathbf{F}\) ）。映射 \(h \mapsto h'\) 是一个同构
\end{enumerate}

\[
\gamma : \operatorname{Hom} _ {\mathbf {C}} (\mathrm{E} \otimes_ {\mathbf {A}} \mathrm{F}, \mathrm{G}) \rightarrow \operatorname{Hom} _ {\mathbf {A}} (\mathrm{E}, \operatorname{Hom} _ {\mathbf {C}} (\mathrm{F}, \mathrm{G})).\tag{6}
\]

特别地，B 和 C 可以取为环 A 的中心的一个子环 \(\Gamma\) ；则对于每个 \(\Gamma\) -模 G，这三个 \(\Gamma\) -模

\[
\operatorname{Hom} _ {\Gamma} (\mathrm{E} \otimes_ {\mathrm{A}} \mathrm{F}, \mathrm{G}), \quad \operatorname{Hom} _ {\mathrm{A}} (\mathrm{E}, \operatorname{Hom} _ {\Gamma} (\mathrm{F}, \mathrm{G})), \quad \operatorname{Hom} _ {\mathrm{A}} (\mathrm{F}, \operatorname{Hom} _ {\Gamma} (\mathrm{E}, \mathrm{G}))
\]

典范同构于 \(\Gamma\) -模，它由 \(\Gamma\) -双线性映射（从 \(\mathbf{E} \times \mathbf{F}\) 到 \(\mathbf{G}\) 中且满足 (1) 的）组成。更具体地说：

推论. 若 C 是交换环，且 E、F、G 是三个 C-模，则 C-模

\[
\begin{array}{l l} \operatorname{Hom} _ {\mathbf {c}} (\mathrm{E} \otimes_ {\mathbf {c}} \mathrm{F}, \mathrm{G}), & \operatorname{Hom} _ {\mathbf {c}} (\mathrm{E}, \operatorname{Hom} _ {\mathbf {c}} (\mathrm{F}, \mathrm{G})), \\ \operatorname{Hom} _ {\mathbf {c}} (\mathrm{F}, \operatorname{Hom} _ {\mathbf {c}} (\mathrm{E}, \mathrm{G})), & \mathcal {L} _ {2} (\mathrm{E}, \mathrm{F}; \mathrm{G}) \end{array}
\]

是典范同构的。
% semantic-fragments: legacy-input-seam
\relax{}
